\documentclass{article}
\usepackage[utf8]{inputenc}
\usepackage{geometry}
\usepackage{enumitem}
\usepackage{hyperref}
\usepackage{xcolor}
\usepackage{graphicx}
\usepackage{array}
\usepackage{multirow}
\usepackage{amsmath, amssymb}
\usepackage{chemfig}
\usepackage{mhchem}
\usepackage{tikz}
\usepackage{setspace}
\usepackage{soul}
\usepackage{mathtools}
\usepackage{booktabs}
\usepackage{chemformula}
\usepackage{siunitx}
\usetikzlibrary{positioning, calc}
\usepackage{longtable}
\geometry{margin=1in}
\setlength{\parskip}{0.5em}
\usepackage{cancel}


\begin{document}

\begin{center}
    \vspace*{-1cm}
    \begin{tabular}{p{12cm} r}
        & \textbf{Name:} \underline{\hspace{5cm}} \\
    \end{tabular}

    \vspace{0.2cm}
    {\LARGE \textbf{Week 10 Worksheet}}
\end{center}    
\\

\vspace{0.2cm}
\noindent\begin{tabular}{|p{0.9\linewidth}|}
\hline
\textbf{(1)} Which of the following solutes is generally soluble in water according to solubility rules?\\[0.2cm]
\begin{minipage}[t]{0.85\linewidth}
\begin{enumerate}[label=\alph*.]
\item BaSO$_4$
\item AgCl
\item KNO$_3$
\item PbI$_2$
\end{enumerate}
\end{minipage}\\[0.2cm]
\hline
\end{tabular}

\noindent\begin{tabular}{|p{0.9\linewidth}|}
\hline
\textbf{(2)} A chemist has 50.0~mL of 2.0~M HCl solution. She adds water to dilute it to a final volume of 200.0~mL. What is the molarity of the diluted solution?\\[0.2cm]
\begin{minipage}[t]{0.85\linewidth}
\begin{enumerate}[label=\alph*.]
\item 0.25~M
\item 0.50~M
\item 1.0~M
\item 2.0~M
\end{enumerate}
\end{minipage}\\[0.2cm]
\hline
\end{tabular}

\noindent\begin{tabular}{|p{0.9\linewidth}|}
\hline
\textbf{(3)} A balloon has a volume of 2.0~L at 300~K. If it is heated to 600~K at constant pressure, what will the new volume be?\\[0.2cm]
\begin{minipage}[t]{0.85\linewidth}
\begin{enumerate}[label=\alph*.]
\item 1.0~L
\item 2.0~L
\item 3.0~L
\item 4.0~L
\end{enumerate}
\end{minipage}\\[0.2cm]
\hline
\end{tabular}

\vspace{0.2cm}
\noindent\begin{tabular}{|p{0.9\linewidth}|}
\hline
\textbf{(4)} A diver's breathing mixture contains helium and oxygen at a total pressure of 8.0~atm. If the partial pressure of oxygen is 1.6~atm, what fraction of the total gas mixture is oxygen?\\[0.2cm]
\begin{minipage}[t]{0.85\linewidth}
\begin{enumerate}[label=\alph*.]
\item 0.10
\item 0.20
\item 0.50
\item 0.80
\end{enumerate}
\end{minipage}\\[0.2cm]
\hline
\end{tabular}
\newpage

\vspace{0.2cm}
\noindent\begin{tabular}{|p{0.9\linewidth}|}
\hline
\textbf{(5)} A 5.0~L sample of hydrogen gas is collected at 310~K. If it is cooled to 155~K at constant pressure, what is the new volume?\\[0.2cm]
\begin{minipage}[t]{0.85\linewidth}
\begin{enumerate}[label=\alph*.]
\item 1.0~L
\item 2.5~L
\item 5.0~L
\item 10.0~L
\end{enumerate}
\end{minipage}\\[0.2cm]
\hline
\end{tabular}

\noindent\begin{tabular}{|p{0.9\linewidth}|}
\hline
\textbf{(6)} A solvent is the substance present in the largest amount in a solution.\\[0.2cm]
\textbf{True} \hspace{1cm} \textbf{False}\\[0.2cm]
\hline
\end{tabular}

\vspace{0.2cm}
\noindent\begin{tabular}{|p{0.9\linewidth}|}
\hline
\textbf{(7)} A sample of gas has a volume of 3.0~L at 27~°C. What will its volume be at 127~°C if pressure remains constant?\\[0.2cm]
\begin{minipage}[t]{0.85\linewidth}
\begin{enumerate}[label=\alph*.]
\item 1.5~L
\item 3.0~L
\item 4.0~L
\item 5.0~L
\end{enumerate}
\end{minipage}\\[0.2cm]
\hline
\end{tabular}

\vspace{0.2cm}
\noindent\begin{tabular}{|p{0.9\linewidth}|}
\hline
\textbf{(8)} A mixture of helium and argon gases in a sealed 10~L container has a total pressure of 5.0~atm. If the partial pressure of helium is 2.0~atm, what is the partial pressure of argon?\\[0.2cm]
\begin{minipage}[t]{0.85\linewidth}
\begin{enumerate}[label=\alph*.]
\item 3.0~atm
\item 7.0~atm
\item 2.5~atm
\item 5.0~atm
\end{enumerate}
\end{minipage}\\[0.2cm]
\hline
\end{tabular}
\newpage

\vspace{0.2cm}
\noindent\begin{tabular}{|p{0.9\linewidth}|}
\hline
\textbf{(9)} A gas sample occupies 400~mL at 273~K. What volume will it occupy at 546~K if pressure remains constant?\\[0.2cm]
\begin{minipage}[t]{0.85\linewidth}
\begin{enumerate}[label=\alph*.]
\item 100~mL
\item 200~mL
\item 400~mL
\item 800~mL
\end{enumerate}
\end{minipage}\\[0.2cm]
\hline
\end{tabular}

\noindent\begin{tabular}{|p{0.9\linewidth}|}
\hline
\textbf{(10)} A drop of water beads up on a freshly waxed car surface. Which statement best explains why this happens?\\[0.2cm]
\begin{minipage}[t]{0.85\linewidth}
\begin{enumerate}[label=\alph*.]
\item Water's adhesion to the wax is stronger than its cohesion.
\item Hydrogen bonding causes strong cohesion within water.
\item London forces between wax molecules attract water.
\item Dipole-induced dipole forces repel water from wax.
\end{enumerate}
\end{minipage}\\[0.2cm]
\hline
\end{tabular}

\vspace{0.2cm}
\noindent\begin{tabular}{|p{0.9\linewidth}|}
\hline
\textbf{(11)} Carlos opens a bottle of perfume and quickly smells it from across the room. Which intermolecular force primarily determines how easily the perfume molecules enter the air?\\[0.2cm]
\begin{minipage}[t]{0.85\linewidth}
\begin{enumerate}[label=\alph*.]
\item Hydrogen bonding
\item Ionic bonding
\item London dispersion forces
\item Metallic bonding
\end{enumerate}
\end{minipage}\\[0.2cm]
\hline
\end{tabular}

\vspace{0.2cm}
\noindent\begin{tabular}{|p{0.9\linewidth}|}
\hline
\textbf{(12)} A student notices that hexane and water do not mix when shaken together. Which type of intermolecular forces explains this behavior?\\[0.2cm]
\begin{minipage}[t]{0.85\linewidth}
\begin{enumerate}[label=\alph*.]
\item London dispersion forces dominate in hexane, but water is polar.
\item Both substances form hydrogen bonds with each other.
\item Hexane is polar, and water is nonpolar.
\item Dipole-dipole interactions attract the two liquids together.
\end{enumerate}
\end{minipage}\\[0.2cm]
\hline
\end{tabular}
\newpage

\vspace{0.2cm}
\noindent\begin{tabular}{|p{0.9\linewidth}|}
\hline
\textbf{(13)} A saturated solution of NaNO$_3$ contains 90~g solute in 100~g water at 50~°C. If 10~g more NaNO$_3$ is added, what happens?\\[0.2cm]
\begin{minipage}[t]{0.85\linewidth}
\begin{enumerate}[label=\alph*.]
\item Dissolves completely
\item Remains undissolved
\item Reacts with water
\item Turns into a supersaturated solution spontaneously
\end{enumerate}
\end{minipage}\\[0.2cm]
\hline
\end{tabular}

\vspace{0.2cm}
\noindent\begin{tabular}{|p{0.9\linewidth}|}
\hline
\textbf{(14)} A weather balloon at 1.20~atm has a volume of 100.~L near the ground. When it rises to a region where the pressure is 0.900~atm, what is its new volume (temperature constant)?\\[0.2cm]
\begin{minipage}[t]{0.85\linewidth}
\begin{enumerate}[label=\alph*.]
\item 75.0~L
\item 90.0~L
\item 100.~L
\item 133~L
\end{enumerate}
\end{minipage}\\[0.2cm]
\hline
\end{tabular}

\vspace{0.2cm}
\noindent\begin{tabular}{|p{0.9\linewidth}|}
\hline
\textbf{(15)} A scuba diver breathes from a tank where the air occupies 12.0~L at 200~atm. If the diver uses half of the air (6.00~L) while underwater at the same temperature, what is the new pressure in the tank?\\[0.2cm]
\begin{minipage}[t]{0.85\linewidth}
\begin{enumerate}[label=\alph*.]
\item 50~atm
\item 100~atm
\item 200~atm
\item 400~atm
\end{enumerate}
\end{minipage}\\[0.2cm]
\hline
\end{tabular}

\vspace{0.2cm}
\noindent\begin{tabular}{|p{0.9\linewidth}|}
\hline
\textbf{(16)} In a sugar-water mixture, sugar is the solvent and water is the solute.\\[0.2cm]
\textbf{True} \hspace{1cm} \textbf{False}\\[0.2cm]
\hline
\end{tabular}
\newpage

\noindent\begin{tabular}{|p{0.9\linewidth}|}
\hline
\textbf{(17)} A diver's air bubble has a volume of 3.0~L at a depth where the pressure is 2.0~atm. If the bubble rises to the surface where the pressure is 1.0~atm, what will its new volume be (assuming constant temperature)?\\[0.2cm]
\begin{minipage}[t]{0.85\linewidth}
\begin{enumerate}[label=\alph*.]
\item 1.5~L
\item 2.0~L
\item 3.0~L
\item 6.0~L
\end{enumerate}
\end{minipage}\\[0.2cm]
\hline
\end{tabular}

\vspace{0.2cm}
\noindent\begin{tabular}{|p{0.9\linewidth}|}
\hline
\textbf{(18)} In percent by volume concentration, both solute and solvent volumes are measured in the same units.\\[0.2cm]
\textbf{True} \hspace{1cm} \textbf{False}\\[0.2cm]
\hline
\end{tabular}

\vspace{0.2cm}
\noindent\begin{tabular}{|p{0.9\linewidth}|}
\hline
\textbf{(19)} When comparing HF and HCl, HF has a much higher boiling point. Which force best explains this difference?\\[0.2cm]
\begin{minipage}[t]{0.85\linewidth}
\begin{enumerate}[label=\alph*.]
\item Dipole-dipole interactions
\item Hydrogen bonding
\item London dispersion forces
\item Covalent bonding
\end{enumerate}
\end{minipage}\\[0.2cm]
\hline
\end{tabular}

\noindent\begin{tabular}{|p{0.9\linewidth}|}
\hline
\textbf{(20)} A piston traps 4.0~L of gas at 1.5~atm. The piston is pushed so that the gas pressure increases to 3.0~atm. What is the new volume of the gas?\\[0.2cm]
\begin{minipage}[t]{0.85\linewidth}
\begin{enumerate}[label=\alph*.]
\item 1.0~L
\item 2.0~L
\item 3.0~L
\item 6.0~L
\end{enumerate}
\end{minipage}\\[0.2cm]
\hline
\end{tabular}
\newpage

\vspace{0.2cm}
\noindent\begin{tabular}{|p{0.9\linewidth}|}
\hline
\textbf{(21)} The concentration of a solution in percent by mass is calculated using the mass of solute divided by the total mass of the solution.\\[0.2cm]
\textbf{True} \hspace{1cm} \textbf{False}\\[0.2cm]
\hline
\end{tabular}

\vspace{0.2cm}
\noindent\begin{tabular}{|p{0.9\linewidth}|}
\hline
\textbf{(22)} A steel tank of fixed volume contains 1.0~mol of nitrogen gas at 2.0~atm. If another 1.0~mol of nitrogen is added at the same temperature, what is the new pressure in the tank?\\[0.2cm]
\begin{minipage}[t]{0.85\linewidth}
\begin{enumerate}[label=\alph*.]
\item 1.0~atm
\item 2.0~atm
\item 3.0~atm
\item 4.0~atm
\end{enumerate}
\end{minipage}\\[0.2cm]
\hline
\end{tabular}

\vspace{0.2cm}
\noindent\begin{tabular}{|p{0.9\linewidth}|}
\hline
\textbf{(23)} A sealed container holds 5.00~L of oxygen gas at 1.00~mol. If half the gas leaks out while temperature and pressure stay constant, what will the new volume be?\\[0.2cm]
\begin{minipage}[t]{0.85\linewidth}
\begin{enumerate}[label=\alph*.]
\item 1.25~L
\item 2.50~L
\item 5.00~L
\item 10.0~L
\end{enumerate}
\end{minipage}\\[0.2cm]
\hline
\end{tabular}

\vspace{0.2cm}
\noindent\begin{tabular}{|p{0.9\linewidth}|}
\hline
\textbf{(24)} What is the molarity of a solution containing 5.00~mol of KBr in 2.00~L of solution?\\[0.2cm]
\begin{minipage}[t]{0.85\linewidth}
\begin{enumerate}[label=\alph*.]
\item 2.00~M
\item 2.50~M
\item 5.00~M
\item 10.0~M
\end{enumerate}
\end{minipage}\\[0.2cm]
\hline
\end{tabular}
\newpage

\vspace{0.2cm}
\noindent\begin{tabular}{|p{0.9\linewidth}|}
\hline
\textbf{(25)} Molarity expresses concentration as moles of solute per liter of solvent.\\[0.2cm]
\textbf{True} \hspace{1cm} \textbf{False}\\[0.2cm]
\hline
\end{tabular}

\vspace{0.2cm}
\noindent\begin{tabular}{|p{0.9\linewidth}|}
\hline
\textbf{(26)} An unsaturated solution contains less solute than the maximum that can dissolve at that temperature.\\[0.2cm]
\textbf{True} \hspace{1cm} \textbf{False}\\[0.2cm]
\hline
\end{tabular}

\noindent\begin{tabular}{|p{0.9\linewidth}|}
\hline
\textbf{(27)} A scuba tank contains 0.80~mol of air in a 6.0~L tank. How many moles of air would fill a 12.0~L tank at the same temperature and pressure?\\[0.2cm]
\begin{minipage}[t]{0.85\linewidth}
\begin{enumerate}[label=\alph*.]
\item 0.400~mol
\item 0.800~mol
\item 1.60~mol
\item 3.20~mol
\end{enumerate}
\end{minipage}\\[0.2cm]
\hline
\end{tabular}

\noindent\begin{tabular}{|p{0.9\linewidth}|}
\hline
\textbf{(28)} A scuba tank contains oxygen at 3.0~atm and nitrogen at 6.0~atm. What is the total pressure in the tank according to Dalton's Law?\\[0.2cm]
\begin{minipage}[t]{0.85\linewidth}
\begin{enumerate}[label=\alph*.]
\item 2.0~atm
\item 3.0~atm
\item 6.0~atm
\item 9.0~atm
\end{enumerate}
\end{minipage}\\[0.2cm]
\hline
\end{tabular}
\newpage

\vspace{0.2cm}
\noindent\begin{tabular}{|p{0.9\linewidth}|}
\hline
\textbf{(29)} According to solubility rules, most silver chloride (AgCl) salts are soluble in water.\\[0.2cm]
\textbf{True} \hspace{1cm} \textbf{False}\\[0.2cm]
\hline
\end{tabular}

\vspace{0.2cm}
\noindent\begin{tabular}{|p{0.9\linewidth}|}
\hline
\textbf{(30)} A gas sample in a steel container has a volume of 12.0~L at 600.~mmHg. If the pressure is reduced to 400.~mmHg, what is the new volume (temperature constant)?\\[0.2cm]
\begin{minipage}[t]{0.85\linewidth}
\begin{enumerate}[label=\alph*.]
\item 8.0~L
\item 12.0~L
\item 18.0~L
\item 20.0~L
\end{enumerate}
\end{minipage}\\[0.2cm]
\hline
\end{tabular}

\vspace{0.2cm}
\noindent\begin{tabular}{|p{0.9\linewidth}|}
\hline
\textbf{(31)} A solution with visible solid particles remaining after stirring is considered supersaturated.\\[0.2cm]
\textbf{True} \hspace{1cm} \textbf{False}\\[0.2cm]
\hline
\end{tabular}

\vspace{0.2cm}
\noindent\begin{tabular}{|p{0.9\linewidth}|}
\hline
\textbf{(32)} A saturated solution can dissolve more solute if the temperature is increased.\\[0.2cm]
\textbf{True} \hspace{1cm} \textbf{False}\\[0.2cm]
\hline
\end{tabular}
\newpage

\vspace{0.2cm}
\noindent\begin{tabular}{|p{0.9\linewidth}|}
\hline
\textbf{(33)} A chemist mixes CO$_2$ and N$_2$ gases in a flask. If CO$_2$ exerts a pressure of 0.750~atm and N$_2$ exerts 0.950~atm, what is the total pressure?\\[0.2cm]
\begin{minipage}[t]{0.85\linewidth}
\begin{enumerate}[label=\alph*.]
\item 0.200~atm
\item 1.70~atm
\item 0.850~atm
\item 1.00~atm
\end{enumerate}
\end{minipage}\\[0.2cm]
\hline
\end{tabular}

\vspace{0.2cm}
\noindent\begin{tabular}{|p{0.9\linewidth}|}
\hline
\textbf{(34)} A balloon inflated to 10.0~L at 293~K is placed in a freezer at 263~K. What is the new volume if pressure is constant?\\[0.2cm]
\begin{minipage}[t]{0.85\linewidth}
\begin{enumerate}[label=\alph*.]
\item 8.00~L
\item 8.98~L
\item 10.0~L
\item 11.0~L
\end{enumerate}
\end{minipage}\\[0.2cm]
\hline
\end{tabular}

\vspace{0.2cm}
\noindent\begin{tabular}{|p{0.9\linewidth}|}
\hline
\textbf{(35)} A 10.0~L rigid container holds 0.50~mol of oxygen gas at 1.5~atm. If the amount of gas doubles while temperature stays constant, what will the final pressure be?\\[0.2cm]
\begin{minipage}[t]{0.85\linewidth}
\begin{enumerate}[label=\alph*.]
\item 0.75~atm
\item 1.5~atm
\item 3.0~atm
\item 6.0~atm
\end{enumerate}
\end{minipage}\\[0.2cm]
\hline
\end{tabular}

\vspace{0.2cm}
\noindent\begin{tabular}{|p{0.9\linewidth}|}
\hline
\textbf{(36)} A weather balloon has a volume of 50.0~L at 250.~K. If it rises and the temperature drops to 200.~K while pressure stays constant, what is the new volume?\\[0.2cm]
\begin{minipage}[t]{0.85\linewidth}
\begin{enumerate}[label=\alph*.]
\item 25.0~L
\item 40.0~L
\item 50.0~L
\item 62.5~L
\end{enumerate}
\end{minipage}\\[0.2cm]
\hline
\end{tabular}
\newpage

\vspace{0.2cm}
\noindent\begin{tabular}{|p{0.9\linewidth}|}
\hline
\textbf{(37)} A tire is inflated with 3.0~mol of air to a volume of 18.0~L. If an additional 1.0~mol of air is added without changing temperature or pressure, what is the final volume?\\[0.2cm]
\begin{minipage}[t]{0.85\linewidth}
\begin{enumerate}[label=\alph*.]
\item 13.5~L
\item 18.0~L
\item 24.0~L
\item 30.0~L
\end{enumerate}
\end{minipage}\\[0.2cm]
\hline
\end{tabular}

\noindent\begin{tabular}{|p{0.9\linewidth}|}
\hline
\textbf{(38)} A balloon contains 1.00~mol of helium and occupies 24.5~L at room conditions. If 2.00~mol of helium are added at the same temperature and pressure, what is the new volume of the balloon?\\[0.2cm]
\begin{minipage}[t]{0.85\linewidth}
\begin{enumerate}[label=\alph*.]
\item 12.3~L
\item 24.5~L
\item 36.8~L
\item 49.0~L
\end{enumerate}
\end{minipage}\\[0.2cm]
\hline
\end{tabular}

\vspace{0.2cm}
\noindent\begin{tabular}{|p{0.9\linewidth}|}
\hline
\textbf{(39)} Two similar compounds, CH$_3$CH$_2$CH$_2$OH and CH$_3$OCH$_2$CH$_3$, have different boiling points. The alcohol boils at a much higher temperature. Why?\\[0.2cm]
\begin{minipage}[t]{0.85\linewidth}
\begin{enumerate}[label=\alph*.]
\item The alcohol has hydrogen bonding, while the ether does not.
\item The ether has stronger dipole-dipole forces.
\item Both compounds have identical London dispersion forces.
\item The ether forms ionic interactions in liquid form.
\end{enumerate}
\end{minipage}\\[0.2cm]
\hline
\end{tabular}

\noindent\begin{tabular}{|p{0.9\linewidth}|}
\hline
\textbf{(40)} Maya accidentally spills water and ethanol on the same counter. She notices that the ethanol evaporates faster. Which statement best explains why?\\[0.2cm]
\begin{minipage}[t]{0.85\linewidth}
\begin{enumerate}[label=\alph*.]
\item Ethanol molecules have stronger hydrogen bonding than water. 
\item Water has stronger hydrogen bonding, which slows evaporation.
\item Ethanol is nonpolar, so it doesn't interact with itself.
\item Water molecules lack dipole-dipole forces.
\end{enumerate}
\end{minipage}\\[0.2cm]
\hline
\end{tabular}
\newpage

\vspace{0.2cm}
\noindent\begin{tabular}{|p{0.9\linewidth}|}
\hline
\textbf{(41)} Two rigid containers are connected by a closed valve. One container holds helium at 2.0~atm in 4.0~L, and the other holds neon at 4.0~atm in 2.0~L. When the valve is opened and gases mix, what is the total pressure in the combined 6.0~L system (assume constant temperature)?\\[0.2cm]
\begin{minipage}[t]{0.85\linewidth}
\begin{enumerate}[label=\alph*.]
\item 2.0~atm
\item 2.7~atm
\item 4.0~atm
\item 6.0~atm
\end{enumerate}
\end{minipage}\\[0.2cm]
\hline
\end{tabular}

\vspace{0.2cm}
\noindent\begin{tabular}{|p{0.9\linewidth}|}
\hline
\textbf{(42)} When cooking pasta, steam rises from the boiling water. Which intermolecular forces are overcome as the water vaporizes?\\[0.2cm]
\begin{minipage}[t]{0.85\linewidth}
\begin{enumerate}[label=\alph*.]
\item Hydrogen bonding between water molecules
\item Ionic attractions between H$^+$ and O$^{2-}$
\item Metallic bonding between hydrogen atoms
\item Covalent bonds within H$_2$O molecules
\end{enumerate}
\end{minipage}\\[0.2cm]
\hline
\end{tabular}

\noindent\begin{tabular}{|p{0.9\linewidth}|}
\hline
\textbf{(43)} A solution is prepared by dissolving 10.0~g of NaCl in 90.0~g of water. What is the mass percent of NaCl in the solution?\\[0.2cm]
\begin{minipage}[t]{0.85\linewidth}
\begin{enumerate}[label=\alph*.]
\item 10.0\%
\item 11.1\%
\item 9.00\%
\item 12.0\%
\end{enumerate}
\end{minipage}\\[0.2cm]
\hline
\end{tabular}

\vspace{0.2cm}
\noindent\begin{tabular}{|p{0.9\linewidth}|}
\hline
\textbf{(44)} A sample of gas occupies 8.0~L at 1.0~atm. The gas is compressed until its volume is 2.0~L. What is the new pressure, assuming temperature is constant?\\[0.2cm]
\begin{minipage}[t]{0.85\linewidth}
\begin{enumerate}[label=\alph*.]
\item 0.25~atm
\item 2.0~atm
\item 4.0~atm
\item 8.0~atm
\end{enumerate}
\end{minipage}\\[0.2cm]
\hline
\end{tabular}
\newpage

\vspace{0.2cm}
\noindent\begin{tabular}{|p{0.9\linewidth}|}
\hline
\textbf{(45)} A balloon containing 2.50~mol of nitrogen occupies 60.0~L. If another 1.50~mol of nitrogen are added at constant temperature and pressure, what is the final volume?\\[0.2cm]
\begin{minipage}[t]{0.85\linewidth}
\begin{enumerate}[label=\alph*.]
\item 36.0~L
\item 60.0~L
\item 96.0~L
\item 120.~L
\end{enumerate}
\end{minipage}\\[0.2cm]
\hline
\end{tabular}

\vspace{0.2cm}
\noindent\begin{tabular}{|p{0.9\linewidth}|}
\hline
\textbf{(46)} According to solubility rules, most nitrate (NO$_3^-$) salts are soluble in water.\\[0.2cm]
\textbf{True} \hspace{1cm} \textbf{False}\\[0.2cm]
\hline
\end{tabular}

\vspace{0.2cm}
\noindent\begin{tabular}{|p{0.9\linewidth}|}
\hline
\textbf{(47)} A lab technician finds that NaCl dissolves easily in water but not in hexane. What is the main reason for this?\\[0.2cm]
\begin{minipage}[t]{0.85\linewidth}
\begin{enumerate}[label=\alph*.]
\item Water's polarity allows it to interact strongly with ionic compounds.
\item Hexane forms hydrogen bonds with Na$^+$ ions.
\item NaCl forms London forces with water.
\item Hexane is polar and repels ions.
\end{enumerate}
\end{minipage}\\[0.2cm]
\hline
\end{tabular}

\vspace{0.2cm}
\noindent\begin{tabular}{|p{0.9\linewidth}|}
\hline
\textbf{(48)} A scientist collects hydrogen gas over water at 25.0~°C. The total pressure is 745~torr, and the vapor pressure of water at 25.0~°C is 24.0~torr. What is the partial pressure of hydrogen?\\[0.2cm]
\begin{minipage}[t]{0.85\linewidth}
\begin{enumerate}[label=\alph*.]
\item 721~torr
\item 769~torr
\item 24.0~torr
\item 745~torr
\end{enumerate}
\end{minipage}\\[0.2cm]
\hline
\end{tabular}
\newpage

\vspace{0.2cm}
\noindent\begin{tabular}{|p{0.9\linewidth}|}
\hline
\textbf{(49)} A gas syringe contains 50.0~mL of nitrogen at 750~torr. If the gas is compressed to 25.0~mL, what is the final pressure (in torr) if temperature remains constant?\\[0.2cm]
\begin{minipage}[t]{0.85\linewidth}
\begin{enumerate}[label=\alph*.]
\item 375~torr
\item 500.~torr
\item $1.0*10^3$~torr
\item $1.5*10^3$~torr
\end{enumerate}
\end{minipage}\\[0.2cm]
\hline
\end{tabular}

\noindent\begin{tabular}{|p{0.9\linewidth}|}
\hline
\textbf{(50)} A balloon is filled with a gas mixture containing 0.6~mol O$_2$ and 0.4~mol N$_2$ at a total pressure of 1.0~atm. What is the partial pressure of O$_2$?\\[0.2cm]
\begin{minipage}[t]{0.85\linewidth}
\begin{enumerate}[label=\alph*.]
\item 0.4~atm
\item 0.6~atm
\item 1.0~atm
\item 1.5~atm
\end{enumerate}
\end{minipage}\\[0.2cm]
\hline
\end{tabular}

\vspace{0.2cm}
\noindent\begin{tabular}{|p{0.9\linewidth}|}
\hline
\textbf{(51)} Dry ice (solid CO$_2$) sublimes directly into a gas at room temperature. Which type of intermolecular force is responsible for holding its molecules together in the solid phase?\\[0.2cm]
\begin{minipage}[t]{0.85\linewidth}
\begin{enumerate}[label=\alph*.]
\item Ionic forces
\item Hydrogen bonding
\item London dispersion forces
\item Dipole-dipole interactions
\end{enumerate}
\end{minipage}\\[0.2cm]
\hline
\end{tabular}

\noindent\begin{tabular}{|p{0.9\linewidth}|}
\hline
\textbf{(52)} Liquid nitrogen boils at -196~°C, while water boils at 100~°C. What difference in intermolecular forces best explains this vast difference?\\[0.2cm]
\begin{minipage}[t]{0.85\linewidth}
\begin{enumerate}[label=\alph*.]
\item Nitrogen has hydrogen bonding; water has dispersion only.
\item Water forms hydrogen bonds; nitrogen has only London forces.
\item Both have identical intermolecular forces.
\item Nitrogen has dipole-dipole interactions stronger than water. 
\end{enumerate}
\end{minipage}\\[0.2cm]
\hline
\end{tabular}
\newpage

\vspace{0.2cm}
\noindent\begin{tabular}{|p{0.9\linewidth}|}
\hline
\textbf{(53)} A sample of hydrogen occupies 10.0~L when it contains 0.400~mol of gas. What would be the volume if the amount were increased to 1.00~mol under the same conditions?\\[0.2cm]
\begin{minipage}[t]{0.85\linewidth}
\begin{enumerate}[label=\alph*.]
\item 4.00~L
\item 10.0~L
\item 25.0~L
\item 40.0~L
\end{enumerate}
\end{minipage}\\[0.2cm]
\hline
\end{tabular}

\vspace{0.2cm}
\noindent\begin{tabular}{|p{0.9\linewidth}|}
\hline
\textbf{(54)} A sealed cylinder of gas has a volume of 8.00~L at 280.~K. If it is heated to 350.~K at constant pressure, what is the final volume?\\[0.2cm]
\begin{minipage}[t]{0.85\linewidth}
\begin{enumerate}[label=\alph*.]
\item 6.40~L
\item 8.00~L
\item 10.0~L
\item 12.0~L
\end{enumerate}
\end{minipage}\\[0.2cm]
\hline
\end{tabular}

\vspace{0.2cm}
\noindent\begin{tabular}{|p{0.9\linewidth}|}
\hline
\textbf{(55)} A rigid 2.0~L cylinder contains 3.0~mol of gas at 5.0~atm. Some gas leaks until only 1.5~mol remain, while temperature and volume remain unchanged. What is the new pressure?\\[0.2cm]
\begin{minipage}[t]{0.85\linewidth}
\begin{enumerate}[label=\alph*.]
\item 1.5~atm
\item 2.5~atm
\item 5.0~atm
\item 10.~atm
\end{enumerate}
\end{minipage}\\[0.2cm]
\hline
\end{tabular}

\vspace{0.2cm}
\noindent\begin{tabular}{|p{0.9\linewidth}|}
\hline
\textbf{(56)} A 15.0~L balloon contains 0.600~mol of helium. If the balloon is filled until it contains 1.50~mol of helium, what will the new volume be at the same temperature and pressure?\\[0.2cm]
\begin{minipage}[t]{0.85\linewidth}
\begin{enumerate}[label=\alph*.]
\item 6.00~L
\item 15.0~L
\item 25.0~L
\item 37.5~L
\end{enumerate}
\end{minipage}\\[0.2cm]
\hline
\end{tabular}
\newpage

\vspace{0.2cm}
\noindent\begin{tabular}{|p{0.9\linewidth}|}
\hline
\textbf{(57)} A solution contains 0.50~mol of NaCl in 250~mL of solution. What is the molarity?\\[0.2cm]
\begin{minipage}[t]{0.85\linewidth}
\begin{enumerate}[label=\alph*.]
\item 0.50~M
\item 1.0~M
\item 2.0~M
\item 0.25~M
\end{enumerate}
\end{minipage}\\[0.2cm]
\hline
\end{tabular}

\vspace{0.2cm}
\noindent\begin{tabular}{|p{0.9\linewidth}|}
\hline
\textbf{(58)} A 10.0~L container holds 0.500~mol of gas at constant temperature and pressure. What volume would 2.00~mol of the same gas occupy under identical conditions?\\[0.2cm]
\begin{minipage}[t]{0.85\linewidth}
\begin{enumerate}[label=\alph*.]
\item 2.50~L
\item 5.00~L
\item 10.0~L
\item 40.0~L
\end{enumerate}
\end{minipage}\\[0.2cm]
\hline
\end{tabular}

\vspace{0.2cm}
\noindent\begin{tabular}{|p{0.9\linewidth}|}
\hline
\textbf{(59)} A chemist dissolves 50.0~mL of ethanol in 150.~mL of water. What is the percent by volume of ethanol?\\[0.2cm]
\begin{minipage}[t]{0.85\linewidth}
\begin{enumerate}[label=\alph*.]
\item 25.0\%
\item 33.3\%
\item 50.0\%
\item 66.7\%
\end{enumerate}
\end{minipage}\\[0.2cm]
\hline
\end{tabular}

\vspace{0.2cm}
\noindent\begin{tabular}{|p{0.9\linewidth}|}
\hline
\textbf{(60)} What is the mass percent of glucose in a solution prepared by dissolving 18~g of glucose in 82~g of water?\\[0.2cm]
\begin{minipage}[t]{0.85\linewidth}
\begin{enumerate}[label=\alph*.]
\item 18\%
\item 20\%
\item 22\%
\item 25\%
\end{enumerate}
\end{minipage}\\[0.2cm]
\hline
\end{tabular}

\end{document}
